\documentclass[UTF8]{ctexart}
\usepackage{xeCJK}
\usepackage{geometry}
\geometry{left=2cm,right=2cm,top=2.5cm,bottom=2.5cm}
\setlength{\headheight}{12.7pt}

\usepackage{titlesec}
\titleformat{\section}
  {\normalfont\large\bfseries}
  {\thesection}
  {1em}
  {}
\titlespacing*{\section}
  {0pt}
  {3.5ex plus 1ex minus .2ex}
  {2.3ex plus .2ex}

\usepackage{amsmath}
\usepackage{amssymb}
\usepackage{amsthm}
\usepackage{enumitem}
\setlist[enumerate]{itemsep=0pt,topsep=0pt,parsep=0pt,leftmargin=0pt,itemindent=2.5em}
\setlist[itemize]{itemsep=0pt,topsep=0pt,parsep=0pt,leftmargin=0pt,itemindent=2.5em}
\usepackage{fancyhdr}
\pagestyle{fancy}
\fancyhf{}
\usepackage{graphicx}
\usepackage{hyperref}
\usepackage{indentfirst}
\usepackage{lmodern}


\begin{document}
\setlength{\abovedisplayskip}{1pt}
\setlength{\belowdisplayskip}{1pt}

\section{子群}

\hypertarget{sec:定义1.1}{}
\noindent\textbf{定义1.1 (子群)}

给定\href{01 半群与群#nameddest=sec:定义2.1}{群} $(G,e,\cdot)$和$H\subset G$.
如果$(H,e,\cdot)$也是群, 就称它是$G$的子群, 记作$H<G$.

\vspace{10pt}

\hypertarget{sec:定理1.1}{}
\noindent\textbf{定理1.1}

任给群$G$和它的非空子集$H$, $H<G$当且仅当对任意的$x,y\in H$有$xy^{-1}\in H$.

\noindent{\kaishu 证明.}

\noindent{\kaishu 必要性.} 显然.

\noindent{\kaishu 充分性.}

任取$x\in H$可得$e=xx^{-1}\in H$; 然后对任意的$x\in H$有$x^{-1}=ex^{-1}\in H$.

最后, 对任意的$x,y\in H$, 有$xy=x(y^{-1})^{-1}\in H$, 即$H$对$\cdot$封闭.
\hfill $\qedsymbol$

\vspace{10pt}

\hypertarget{sec:定理1.2}{}
\noindent\textbf{定理1.2}

任给群$G$和它的一族子集$\mathcal{H}$, 它们的交集$\displaystyle\bigcap\mathcal{H}$
也是$G$的子集.

\noindent{\kaishu 证明.}

如果$\displaystyle x,y\in\bigcap\mathcal{H}$, 那么
对任意的$H\in\mathcal{H}$都有$x,y\in H\implies xy^{-1}\in H$,
所以$\displaystyle xy^{-1}\in\bigcap\mathcal{H}$. \hfill $\qedsymbol$

\vspace{10pt}

\hypertarget{sec:定义1.2}{}
\noindent\textbf{定义1.2 (生成子群)}

给定群$G$和它的子集$S$, 定义
\[
    \langle S\rangle\overset{\mathrm{def}}{=}
    \bigcap\,\{H\subset G\mid H<G\,\land\,H\supset S\},\\[1pt]
\]
它是$G$的包含$S$的最小子群, 称之为$S$在$G$中的\textbf{生成子群}.

\noindent{\kaishu 验证.}

第一, 记$\mathcal{H}=\{H\subset G\mid H<G\,\land\,H\supset S\}$,
它的元素都是$G$的子群,
从而$\displaystyle\langle S\rangle=\bigcap\mathcal{H}$是$G$的子群.

第二, $\mathcal{H}$的元素都包含$S$,
从而$\displaystyle\langle S\rangle=\bigcap\mathcal{H}$包含$S$.

第三, 如果$H$也是$G$的包含$S$的子群, 那么$H\in\mathcal{H}$,
从而$\displaystyle\langle S\rangle=\bigcap\mathcal{H}$包含于$H$.
\hfill $\qedsymbol$

\vspace{10pt}

\hypertarget{sec:定理1.3}{}
\noindent\textbf{定理1.3}

任给群$G$, 它的所有子群在包含关系下构成一个格, 其中对任意的$H,K<G$,
\[
    H\lor K=\langle H\cup K\rangle,\quad H\land K=H\cap K.
\]

\noindent{\kaishu 证明.}

第一, 显然$\langle H\cup K\rangle$包含$H,K$, 并且如果$L$也包含$H,K$, 则
\[
    H\cup K\subset L\quad\implies\quad\langle H\cup K\rangle\subset L,
\]
第二, 显然$H\cap K$包含于$H,K$, 并且如果$L$也包含于$H,K$, 则$L\subset H\cap K$.
\hfill $\qedsymbol$

\end{document}